% SPDX-License-Identifier: Apache-2.0
\section{Results}
\label{sec:results}

On \ReportDate, \NRtlPass{} of the \NTests{} tests passed on Vreteno's
netlist. The failures, \NRtlFail{} in all, divide into two kinds:
known failures, each tracked by an issue (\NKnownFail), and new ones
(\NUnknownFail).
\NSailPass{} of the \NTests{} tests passed their self-check on Sail. The
Sail column is a check on the harness rather than on the core: each
self-checking program must pass on the model that wrote its expected
values. When it does, the description of the core, the compiler flags
and the test sources agree with one another.

Table~\ref{tab:suites} gives the results by suite. The checks column
counts the values a test compares with Sail's. The counts exclude the
fill of the unused signature area and the marker words. The suite that
compares the most is \code{\MostChecksSuite}, with \MostChecks{}
values in \MostChecksTests{} tests.

\begin{table}[t]
\centering
\caption{Results by suite}
\label{tab:suites}
\scriptsize
\input{suites}
\end{table}

\subsection{Cycles}

The \NTests{} runs took \TotalCycles{} cycles to retire \TotalInstr{}
instructions, \MeanCPI{} cycles per instruction on average.
Fig.~\ref{fig:cpi} breaks this down by suite.

These numbers describe the testbench as much as the core. An
architectural test is long straight-line code that runs once. The
instruction cache then helps only within a line of four words, and
each new line is fetched over the bus. The slowest suite is \code{\SlowSuite}, at \SlowCPI{} cycles
per instruction, and the fastest is \code{\FastSuite}, at \FastCPI.
The multiply and divide instructions take several cycles each in
Vreteno's sequencer, which is why \code{M} and \code{Zmmul} are above
the rest. The figures are a record that the runs completed
at a plausible speed. They are not a performance claim.

\begin{figure}[t]
\centering
\input{cpichart}
\caption{Mean cycles per instruction on the netlist, by suite.}
\label{fig:cpi}
\end{figure}
